\subsection{A second design that obstructs geometric localization}
\label{geo:singer}

The first-moment proof used a coupling on well-conditioned projective
frames. A natural stronger localization at second moment would couple
nine points whose blowup has nef anticanonical class. Here \emph{nef}
means nonnegative intersection with every curve. The next example
shows that second-moment isotropy alone cannot enforce that geometry.

Let \(\zeta=e^{2\pi i/7}\), and set
\begin{equation}
 B=\{e_0,e_1,e_2\},\qquad
 T_j=\frac{(1,\zeta^j,\zeta^{3j})}{\sqrt3},\quad 0\le j<7,\qquad
 \mu_{\rm S}=\frac1{12}\sum_{b\in B}\delta_b+
                  \frac3{28}\sum_{j=0}^6\delta_{T_j}.
 \label{geo:singer-measure}
\end{equation}
This ten-point design is a known projective-toric construction
\cite[Example 4.2]{geo_iosue2024}. Its character exponents
\(\{0,2,6,1,3,4\}\) in the symmetric-square basis are distinct modulo
seven, giving \(M_2(\mu_{\rm S})=I_6/6\) by exact character cancellation.

\begin{theorem}[An exact coupling obstruction and repair threshold]
\label{geo:singer-nef}
There is no coupling with all marginals \(\mu_{\rm S}\) supported on
nine distinct points whose blowup of \(\mathbb P^2\) has nef
anticanonical class. For
\[
 \mu_t=t\mu_{\rm S}+(1-t)\sigma,
\]
such a coupling exists if and only if \(t\le20/21\).
The coupling may always be taken exchangeable.
\end{theorem}

\begin{proof}
The cubics
\[
 F_0=x^3-yz^2,\qquad F_1=y^3-x^2z,\qquad F_2=z^3-xy^2
\]
vanish at all \(T_j\). Each omits one basis point and has a cusp at
another. A nine-point subset obtained by deleting a basis point thus
has a cubic through eight points simply and one doubly; its strict
transform has anticanonical intersection \(9-10=-1\).
Deleting a torus point instead leaves nine points on a smooth cubic.
For deletion of \(T_0\), an example is
\[
 x^2y+x^2z+xy^2+y^2z+xz^2+yz^2+xyz=0.
\]
Its gradient differences are
\[
 G_x-G_y=(y-x)(x+y-z),\quad
 G_y-G_z=(z-y)(y+z-x),\quad
 G_z-G_x=(x-z)(z+x-y).
\]
At a common gradient zero, if \(x=y=z\), a gradient entry is \(7x^2\).
If exactly two coordinates agree, say \(x=y\ne z\), the remaining
differences force \(z=0\), but then \(G_z=3x^2\ne0\).
If all three are distinct, the differences force
\(x+y=z,\ y+z=x,\ z+x=y\), hence the zero vector.
The cubic is therefore smooth, and
the diagonal order-seven symmetry handles all other deletions.
A smooth cubic through nine points has a strict transform of square
zero; its intersections with every distinct irreducible curve are
nonnegative, so it is nef.

Thus a nef nine-subset of the ten-point support contains all three
basis points. Its expected number of basis points would be three,
whereas its prescribed marginals give \(9\mu_{\rm S}(B)=9/4\).

For a tuple with marginal \(\mu_t\), let \(n_T\) count entries in the
seven-point torus orbit, and \(n_O\) entries outside the ten-point
support. Every admissible tuple satisfies \(n_T\le6+n_O\): when
\(n_O=0\) it must omit a torus point, and otherwise distinctness gives
the assertion. Taking expectations yields
\[
 \frac{27t}{4}\le6+9(1-t),\qquad t\le\frac{20}{21}.
\]
At equality, mix two types and randomly permute the tuple:
with probability \(4/7\), all three basis points and six uniformly
chosen torus points; with probability \(3/7\), all seven torus points,
one uniform basis point, and one Haar point.
For the second type, cyclic symmetry reduces to the retained basis
point \(e_0\). Cubics through \(T\cup\{e_0\}\) form the pencil
\(\operatorname{span}(F_1,F_2)\), since the seven cubic evaluation
characters exhaust all residues modulo seven.
This pencil contains the smooth member
\[
 F_1+F_2=y^3-x^2z+z^3-xy^2.
\]
Indeed, its \(y\)-derivative is \(y(3y-2x)\). If \(y=0\), the other
gradient equations force the zero vector. Otherwise \(y=2x/3\),
the \(x\)-derivative forces \(z=-2x/9\), and the \(z\)-derivative is
\((4/27-1)x^2\ne0\). Only finitely many pencil members are singular,
so a Haar point chooses a smooth member almost surely.
The mixture has weights \(5/63\) at each basis point, \(5/49\) at
each torus point, and Haar weight \(1/21\), exactly \(\mu_{20/21}\).
For smaller \(t\), mix with independent Haar points, whose generic
nine-point configuration lies on a smooth cubic.
\end{proof}

Two measures with identical second moments can therefore have
different coupling geometry. The obstruction is not explained by a
large atom: every atom of \(\mu_{\rm S}\) has mass below \(1/9\).
It is encoded in the incidence and multiplicities of the three cubics.

\subsection{The same design produces a singular K3 surface}
\label{geo:k3}

Blow up the ten points and call the rational surface \(S\). The strict
transforms \(C_0,C_1,C_2\) of the preceding cubics are disjoint smooth
rational curves with
\[
 C_i^2=-3,\qquad C_iC_j=0\ (i\ne j),\qquad
 C_0+C_1+C_2=-3K_S,\qquad K_S^2=-1.
\]
The cyclic triple cover \(\pi:Z\to S\), using the line bundle
\(-K_S\) and branch divisor \(\sum C_i\), is smooth: locally it has
equation \(w^3=u\). Its ramification curves \(R_i\) have square
\(-1\), and the canonical bundle formula gives
\(K_Z=R_0+R_1+R_2\). Contract them to obtain \(X\).
Then \(K_X\) is trivial, and
\[
 e(Z)=3e(S)-2\sum_ie(C_i)=3\cdot13-2\cdot6=27,\qquad e(X)=24.
\]
Noether's formula gives \(\chi(\mathcal O_X)=2\); connectedness and
triviality of \(K_X\) imply \(h^1(\mathcal O_X)=0\).
Thus \(X\) is a smooth projective K3 surface.
The term \emph{singular K3} below means maximal Picard rank twenty.

\fig{fig05_singer_cover}{From a finite projective design to a K3 surface.
(a) The three coordinate rays and seven Singer rays are drawn as an
incidence schematic; their positions in the drawing are not a projective
embedding. (b) Blowing up the ten rays separates the three cubic transforms
into disjoint $(-3)$-curves. The triple cover makes their ramification
curves $(-1)$-curves; contracting these gives a smooth K3 surface with
Euler number 24. (c) The transcendental lattice has the Eisenstein shape
$A_2(7)$, Gram matrix $\left(\begin{smallmatrix}14&-7\\-7&14\end{smallmatrix}\right)$,
and discriminant 147. The arrows in (b) describe actual geometric
operations, while (c) is a lattice diagram.}{fig:k3}

\begin{theorem}[The lattice of the design cover]
\label{geo:k3-lattice}
The K3 surface obtained from this ten-point configuration satisfies
\[
 \rho(X)=20,\qquad
 T_X\cong A_2(7)=
 \begin{pmatrix}14&-7\\-7&14\end{pmatrix},\qquad
 \operatorname{disc}T_X=147.
\]
It admits a symplectic automorphism group \(C_7\rtimes C_3\) and a
commuting non-symplectic deck transformation of order three with
exactly three isolated fixed points.
\end{theorem}

\begin{proof}
A weighted-projective model is
\[
 w^3=(x^3-yz^2)(y^3-x^2z)(z^3-xy^2),\qquad
 \deg(x,y,z,w)=(1,1,1,3).
\]
On its smooth locus a two-form is
\[
 \Omega=\frac{x\,dy\wedge dz-y\,dx\wedge dz+z\,dx\wedge dy}{w}.
\]
At each blown-up point the total branch multiplicity is three, so
the blowup removes its apparent poles. The zeros are exactly the
ramification curves, removed by their contraction.
The deck map \(\delta:w\mapsto e^{2\pi i/3}w\) acts nontrivially on
\(\Omega\); its three fixed curves become three isolated fixed points.
The maps
\[
 g:(x,y,z,w)\mapsto(x,\zeta y,\zeta^3z,\zeta^4w),\qquad
 c:(x,y,z,w)\mapsto(y,z,x,w)
\]
preserve \(\Omega\), and satisfy \(g^7=c^3=1\),
\(cgc^{-1}=g^2\) in weighted projective coordinates.

A symplectic order-seven automorphism has rank-four invariant
cohomology and a rank-eighteen algebraic coinvariant lattice
\cite[Propositions 4.5--4.6]{geo_garbagnati_sarti2007}.
An invariant ample class therefore gives \(\rho(X)\ge19\).
The non-symplectic order-three action makes the transcendental
lattice a module over \(\mathbb Z[e^{2\pi i/3}]\), so its rank is
even. It is at most three and at least two, hence is two.
Thus \(\rho(X)=20\). Since the Eisenstein ring is principal, its
invariant positive even rank-two form is \(A_2(m)\) for an integer \(m>0\).

To determine \(m\), let
\[
 L_0=\{D\in\operatorname{Pic}(S):D C_i=0\text{ for every }i\}.
\]
The primitive closure of \(\langle C_0,C_1,C_2\rangle\) has index
three: it contains \(K_S=-(C_0+C_1+C_2)/3\), and its index squared
divides \(27\). Unimodularity of \(\operatorname{Pic}(S)\) gives
\(|\operatorname{disc}L_0|=3\).
Pullback scales its form by three. The fixed lattice of \(\delta\)
has rank eight and discriminant \(3^7\), since its fixed locus is
three isolated points \cite[Theorem 4.1]{geo_artebani_sarti_taki2011}.
Hence \(\pi^*L_0\) has index three in that fixed lattice.

Let \(H\) and \(E_b,E_{T_j}\) be the line and exceptional classes on \(S\).
The \(g\)-fixed part of \(L_0\) has basis
\[
 v_1=H-\sum_{b\in B}E_b,\qquad
 v_2=7H-3\sum_{j=0}^6E_{T_j},\qquad
 \bigl(v_iv_j\bigr)=\begin{pmatrix}-2&7\\7&-14\end{pmatrix}.
\]
The order-seven action on the index-three quotient is trivial, and
the obstruction in \(H^1(C_7,\pi^*L_0)\) is seven-torsion. Thus
taking \(g\)-invariants preserves index three.
The unique even index-three overlattice of the threefold-scaled
matrix adds \((v_2-v_1)/3\): an order-three discriminant element
\((\alpha v_1+\beta v_2)/3\) is isotropic exactly when
\(\alpha+\beta=0\) modulo three.
In the basis \(v_1+(v_2-v_1)/3,-(v_2-v_1)/3\) its matrix is
\[
 N=\begin{pmatrix}2&1\\1&-10\end{pmatrix}.
\]
A symplectic automorphism fixes \(T_X\) pointwise, since its rational
coinvariant classes are of type \((1,1)\).
The rank-four invariant lattice therefore has rank-two algebraic part.
This contains the rank-two lattice just computed in the \(\delta\)-fixed
part; consequently \(\delta\) acts trivially on its rational span, and
\(\operatorname{NS}(X)^g=N\).

The full invariant lattice is
\[
 H^2(X,\mathbb Z)^g
 \cong U(7)\oplus\begin{pmatrix}4&1\\1&2\end{pmatrix},
 \qquad |\operatorname{disc}|=7^3,
\]
with elementary discriminant group \((\mathbb Z/7)^3\)
\cite[Proposition 4.6]{geo_garbagnati_sarti2007}.
Both \(N\) and \(T_X\) are primitive in this lattice.
If \(h\) is their gluing index, the gluing group injects into both
discriminant groups, so \(h\mid21\). Determinants give
\[
 7^3=\frac{21\cdot3m^2}{h^2}.
\]
The only possibilities are \((m,h)=(7,3),(49,21)\).
In the latter case, the seven-primary discriminants of \(N,T_X\)
are \(\mathbb Z/7\) and \((\mathbb Z/49)^2\). The gluing subgroup
has generator \((n,7v)\) with \(n\ne0\).
Its orthogonal complement projects onto \((\mathbb Z/49)^2\),
since every pairing with \(7v\) can be canceled by a multiple of \(n\).
Lift an element of order \(49\). Seven times that lift has zero first
component and nonzero second component; it cannot belong to the
gluing subgroup. The glued discriminant would retain an element of
order \(49\), contradicting its elementary form. Therefore \(m=7\).
\end{proof}

The invariant lattice also makes the relevant nef geometry explicit.
Its norm is \(2(x^2+xy-5y^2)\), congruent to \(2(x-y)^2\) modulo three,
so it represents neither \(-2\) nor \(-14\).
A rational curve on \(X\) has square \(-2\); its \(g\)-orbit cannot
have length one. A length-seven orbit sum has square \(-14+14k\)
for an integer \(k\ge0\), and cannot have square \(-14\), hence is
nonnegative. Classes in the positive component of \(N_{\mathbb R}\)
therefore pair positively with every such orbit sum, and thus with
every curve, by the Hodge index theorem. The invariant nef cone is
the closure of that positive component. In invariant blowup
coordinates
\[
 dH-b\sum_{b\in B}E_b-c\sum_jE_{T_j},\qquad
 3d=3b+7c,
\]
its null boundary rays are \(d/b=(9\pm\sqrt{21})/2\).

The construction does not assert a new K3 isomorphism class. Its
content is the explicit passage from an isotropic quantum ensemble,
through a geometric localization obstruction, to a canonical cover
whose transcendental lattice is determined. The same cubic
multiplicities that obstruct nef localization make the cyclic cover
possible.

\subsection{Verification and reproducible mathematical content}
\label{geo:verification}

The conceptual inputs are Kempf--Ness, the three-qutrit Cartan
classification, semicontinuity of singularity exponents,
Hilbert/Bergman geodesic comparison, and the stated K3 automorphism
lattices. They enter with precise hypotheses and primary references.
Finite calculations are supplied as exact-arithmetic certificates:
Moore determinants and quartic contractions; Hesse moments, Gram
matrices and spectra; Singer character cancellation, cubic
smoothness and intersections; and finite minimal-jet moment-map
examples. Numerical sampling illustrates bounds and searches for
counterexamples. The theorem constants follow from the stated
inequalities and identities.

The operational supplement evaluates \eqref{geo:confidence} from
measurement counts and supplies seeded synthetic-data examples.
It exposes the equation-to-projector and moment-to-weight-conditioning
steps of \eqref{geo:general-cert} for reuse with other finite ray
sets. The computations verify finite identities and statistical
calculations; they are not a completed machine formalization of
the analytic and geometric proofs.
